\chapter{CONCLUSÃO}
\label{chapter:conclusao}

Este trabalho concluiu o ciclo de projeto mecânico previsto na disciplina de Dinâmica: a fundamentação teórica, a construção do protótipo, o levantamento dos parâmetros físicos e o equacionamento completo do movimento do robô articulado que se desloca sobre uma corda.

A elaboração da modelagem teórica e física do mecanismo articulado demonstrou a importância da padronização dos sistemas de referência e do uso rigoroso das equações vetoriais na engenharia. A utilização das equações de velocidade e aceleração relativa permitiu transcrever a geometria do movimento do robô para uma formulação matemática consistente.

A etapa de levantamento dos parâmetros físicos do protótipo construído com palitos de madeira e motor DC — extraindo massas, dimensões e calculando tensores de inércia por decomposição em sólidos homogêneos e composição pelo teorema dos eixos paralelos — constituiu a base fundamental para a aplicação das leis de Newton e da equação de Euler para a cinética do sistema. O robô montado resultou em 103,1 g, com centro de massa 21,4 mm abaixo do eixo do motor, valores dominados pela bateria e pelo motor com caixa de redução, que somam 91\,\% da massa.

O modelo cinemático e dinâmico, construído pela cadeia de referenciais e pelas equações de Newton-Euler projetadas nas velocidades generalizadas, foi verificado contra a formulação de Lagrange com diferença relativa da ordem de $10^{-16}$. A simulação com o motor a 30 rpm no sentido anti-horário previu um avanço para a frente do centro de massa de 130,8 mm em 8 s (velocidade média de 22,6 mm/s, cerca de 52\,\% do limite da cinemática pura de 43,5 mm/s), com força vertical da corda sempre positiva e trocas de mão por impacto plástico. O mecanismo do avanço — engate do gancho à frente, seguido do balanço pendular do corpo — ficou claramente caracterizado nas trajetórias e nos degraus da curva de deslocamento.

Espera-se que, com o equacionamento matricial e vetorial estabelecido neste relatório, a futura comparação com as filmagens do protótipo valide o deslocamento previsto, e que o cálculo do torque necessário no eixo do motor — obtido como força generalizada associada à rotação da manivela — permita dimensionar o motor e as reações nas juntas de madeira durante a operação, completando o detalhamento dinâmico do projeto.

As principais extensões do trabalho são: medir a massa do motor e o amortecimento real do mecanismo; derivar a regra de engate da geometria precisa do gancho; incluir a rolagem em torno da corda em um modelo espacial; e comparar os ângulos do chassi e das pernas previstos com os medidos nas filmagens do protótipo em operação.
